\scanpage{12}
D'autre part, on a une résolution injective de $\underline O_X$ par
\[
 \underline O_X\longrightarrow\underbrace{\underline K_X\longrightarrow\underline K_X/\underline O_X}_{K^\bullet}\longrightarrow0\longrightarrow0\longrightarrow0\ldots
\]
d'où
\[
 \operatorname{Ext}^{*}(X;F,\mathcal O_X)\simeq H^{*}(0\longrightarrow
 \underset{0}{\operatorname{Hom}(M_1,K)}\longrightarrow\underset{1}{\operatorname{Hom}(M_0,K/A)}
 \longrightarrow0\longrightarrow\ldots)
\]
i.e. nul en dim $\ne0,1$ et
\[
\begin{aligned}
 \unsure{\overline{\mathbb E}|\mathbb E|E}\mathrm{xt}^0(X;F,\mathcal O_X)&\simeq\operatorname{Ker}(\operatorname{Hom}(M_1,K)\longrightarrow\operatorname{Hom}(M_0,K/A))\\
 \operatorname{Ext}^1(X;F,\mathcal O_X)&\simeq\operatorname{Coker}(\qquad\mathrm{id}\qquad)
\end{aligned}
\]
Donc le th. de dualité devient
\[
\begin{aligned}
 E^0&\xrightarrow[\sim]{\varphi^0}\operatorname{Hom}(\operatorname{Coker}u,K/A)\\
 E^1&\xrightarrow[\sim]{\varphi^1}\operatorname{Hom}(\operatorname{Ker}u,K/A)
\end{aligned}
\]
qui résulte des faits suivants

a) Tout hom $M_1\longrightarrow A$ est nul (l'image est un sous-groupe $K$-vectoriel de $A$) donc $\varphi^0$ injectif

b) Tout hom. $M_1\longrightarrow K/A$ se relève en $M_{\unsure{1|a}}\longrightarrow K$, i.e. $\operatorname{Ext}^1_A(M_1,A)=0$, ou aussi \uline{$\operatorname{Ext}^1_A(K,A)=0$} [Pour le voir, utiliser la résolution injective \deleted{$A\to$} de $A$ par le complexe $0\longrightarrow K\longrightarrow K/A\longrightarrow0$, on obtient \deleted{\unsure{le}} complexe $0\longrightarrow\underset{\substack{\|\\K}}{\operatorname{Hom}_A(K,K)}\longrightarrow\underset{\substack{\unsure{\wr}\\\widehat K=K}}{\operatorname{Hom}(K,K/A)}\longrightarrow0$, qui est acyclique - utiliser \uline{autodualité} de $K$ \unsure{par} \inserted{par \deleted{\unsure{le}} le fait que $A$ \uline{complet}} ]. Donc $\varphi^0$ surjectif

c) $K/A$ est injectif, donc tt hom $\operatorname{Ker}u\longrightarrow K/A$ s'étend : $M_0\longrightarrow K/A$, donc $\varphi^1$ est surjectif, et tout hom de $M_0/\operatorname{Ker}u\subset M_1$ dans $K/A$ s'étend en un hom $M_1\longrightarrow K/A$, qui on vient de voir se remonte en $M_1\longrightarrow K$, donc $\varphi^1$ est injectif.
% archive sheet number: 11

